\documentclass[a4paper, serif, 10pt]{moderncv}

%% moderncv
% style and color
\moderncvstyle{banking}
\moderncvcolor{black}

% renew moderncv command to adapt for CJK colon
\renewcommand*{\cvitem}[3][.25em]{%
  \ifstrempty{#2}{}{\hintstyle{#2}：}{#3}%
  \par\addvspace{#1}}

\renewcommand*{\cvitemwithcomment}[4][.25em]{%
  \savebox{\cvitemwithcommentmainbox}{\ifstrempty{#2}{}{\hintstyle{#2}：}#3}%
  \setlength{\cvitemwithcommentmainlength}{\widthof{\usebox{\cvitemwithcommentmainbox}}}%
  \setlength{\cvitemwithcommentcommentlength}{\maincolumnwidth-\separatorcolumnwidth-\cvitemwithcommentmainlength}%
  \begin{minipage}[t]{\cvitemwithcommentmainlength}\usebox{\cvitemwithcommentmainbox}\end{minipage}%
  \hfill% fill of \separatorcolumnwidth
  \begin{minipage}[t]{\cvitemwithcommentcommentlength}\raggedleft\small\itshape#4\end{minipage}%
  \par\addvspace{#1}}

%% page layout/margins
\usepackage[top=2.5cm, bottom=2.5cm, left=1.5cm, right=1.5cm]{geometry}

\nopagenumbers{}

%% Babel config for English language
\usepackage[english]{babel}

%% fontspec
\usepackage{fontspec}

\IfFontExistsTF{Linux Libertine}{
  \setmainfont[Ligatures={TeX, Common}, Numbers=Lining, ItalicFont=Linux Libertine]{Linux Libertine}
}{}
\IfFontExistsTF{Linux Libertine O}{
  \setmainfont[Ligatures={TeX, Common}, Numbers=Lining, ItalicFont=Linux Libertine O]{Linux Libertine O}
}{}

%% CTeX
% CJK support, used to show CJK characters in the resume
%
% - fontset=none: disable builtin fontset but instead set the CJK font manually
% - heading=false: disable ctex heading
% - punct=kaiming: use kaiming punctuations styles for CJK
% - scheme=plain: use plain scheme, do not override \normalsize font size
% - space=auto: space settings for CJK characters
%
% ref:
% - http://ctan.mirrorcatalogs.com/language/chinese/ctex/ctex.pdf
\usepackage[UTF8, heading=false, punct=kaiming, scheme=plain, space=auto]{ctex}

\IfFontExistsTF{Noto Serif CJK SC}{
  \setCJKmainfont{Noto Serif CJK SC}
}{}
\IfFontExistsTF{Noto Sans CJK SC}{
  \setCJKsansfont{Noto Sans CJK SC}
}{}

%% line spacing
\usepackage{setspace}
\setstretch{1.125}

%% URL styling - use normal text instead of monospace
\urlstyle{same}

% Auto-underline all links
% Defer until \begin{document} so hyperref is already loaded
\AtBeginDocument{%
  \let\oldhref\href
  \renewcommand{\href}[2]{\oldhref{#1}{\underline{#2}}}%
}

%% Patch \httplink and \httpslink to handle full URLs
% The original moderncv macros always prepend http:// or https:// to URLs.
% This causes issues when the URL already has a protocol (e.g., https://example.com)
% resulting in malformed URLs like https://https://example.com.
% This patch checks if the URL already contains "://" and uses it directly if so.
% Note: We use \str_set:Nx (x-expansion) to expand macros like \@homepage first.
\ExplSyntaxOn
\str_new:N \g_yamlresume_url_str

\RenewDocumentCommand{\httpslink}{O{}m}{
  \str_set:Nx \g_yamlresume_url_str {#2}
  \str_if_in:NnTF \g_yamlresume_url_str {://}
    { \url{#2} }
    { \url{https://#2} }
}

\RenewDocumentCommand{\httplink}{O{}m}{
  \str_set:Nx \g_yamlresume_url_str {#2}
  \str_if_in:NnTF \g_yamlresume_url_str {://}
    { \url{#2} }
    { \url{http://#2} }
}
\ExplSyntaxOff

\name{陳詠欣}{}
\title{數據科學家｜機器學習與大數據分析專家}
\phone[mobile]{+852 9876 5432}
\email{wingyan@wingyanchan.com}
\homepage{https://www.wingyanchan.com}

\address{中國香港，香港島，香港，香港中環皇后大道中100號}{}{}

\extrainfo{{\small \faGithub}\ \href{https://github.com/wingyanchan}{@wingyanchan} {} {} {} • {} {} {}
{\small \faLinkedin}\ \href{https://www.linkedin.com/in/wingyanchan}{@wingyanchan} {} {} {} • {} {} {}
{\small \faMedium}\ \href{https://medium.com/@wingyanchan}{@wingyanchan}}

\begin{document}

\maketitle

\section{簡介}

\cvline{}{擁有超過 6 年數據科學及機器學習經驗，專注於金融科技、智慧城市及零售分析等領域。
%
\begin{itemize}
\item 善於運用統計建模、機器學習及深度學習技術解決實際商業問題
\item 熟悉大數據平台（Spark、Hadoop）及雲端服務（AWS、Azure）
\item 具備端到端數據產品開發經驗，由數據收集、特徵工程到模型部署及監控
\item 熱衷於把複雜的數據分析結果轉化為清晰的商業決策建議
\item 擁有多項專業認證及期刊論文發表經驗
\end{itemize}}
\section{教育背景}

\cventry{2017年9月–2019年7月}
        {碩士，數據科學，成績：3.8/4.3}
        {香港大學}
        {\url{https://www.hku.hk/}}
        {}
        {\begin{itemize}
\item 以優異成績（Distinction）畢業，專注於機器學習及大數據分析
\item 畢業論文題目為「基於深度學習的香港樓價預測模型」，獲學院最佳論文獎
\item 擔任研究助理，協助教授進行自然語言處理相關研究
\end{itemize}
\textbf{課程}：高級機器學習、大數據分析、統計學習理論、深度學習、數據可視化、雲端運算}

\cventry{2013年9月–2017年7月}
        {學士，計算機科學，成績：3.5/4.3}
        {香港中文大學}
        {\url{https://www.cuhk.edu.hk/}}
        {}
        {\begin{itemize}
\item 副修統計學，建立扎實的數學及程式設計基礎
\item 參與多個編程比賽及黑客松，曾獲校內數據分析比賽冠軍
\end{itemize}
\textbf{課程}：數據結構、算法設計、數據庫系統、概率與統計、人工智能}
\section{工作經歷}

\cventry{2020年1月至今}
        {高級數據科學家}
        {智庫數據科技有限公司}
        {\url{https://www.thinktankdata.io}}
        {}
        {\begin{itemize}
\item 領導數據科學團隊，為銀行及零售客戶開發風險評估及客戶分群模型
\item 設計並部署基於 Python 及 Spark 的機器學習平台，將模型訓練時間縮短 40\%
\item 與業務部門緊密合作，將分析結果轉化為可執行的商業策略
\item 推動 MLOps 實踐，建立自動化模型監控及重新訓練流程
\item 指導初級數據科學家，並定期舉辦內部培訓
\end{itemize}
\textbf{關鍵字}：Python、Spark、MLOps、機器學習、風險建模}

\cventry{2019年8月–2019年12月}
        {數據科學家}
        {香港電訊有限公司}
        {\url{https://www.hkt.com}}
        {}
        {\begin{itemize}
\item 分析流動通訊客戶的使用行為及流失率，建立客戶流失預警模型
\item 利用 SQL 及 Hive 處理每日數十億條數據記錄，製作儀表板供管理層參考
\item 開發客戶生命週期價值（LTV）模型，協助市場部制定精準營銷策略
\end{itemize}
\textbf{關鍵字}：SQL、Hive、客戶流失、Tableau}
\section{語言}

\cvline{粵語}{母語或雙語水平 \hfill \textbf{關鍵字}：廣東話}
\cvline{英語}{完全專業水平 \hfill \textbf{關鍵字}：IELTS 8.0}
\cvline{普通話}{完全專業水平 \hfill \textbf{關鍵字}：普通話水平測試二級甲等}
\section{專業技能}

\cvline{機器學習}{專家 \hfill \textbf{關鍵字}：監督式學習、非監督式學習、深度學習、模型評估、特徵工程、強化學習}
\cvline{程式語言及工具}{專家 \hfill \textbf{關鍵字}：Python、R、SQL、Java、TensorFlow、PyTorch、scikit-learn}
\cvline{大數據平台}{高級 \hfill \textbf{關鍵字}：Apache Spark、Hadoop、Hive、Kafka、Airflow、Databricks}
\cvline{雲端及 MLOps}{高級 \hfill \textbf{關鍵字}：AWS、Azure、Docker、Kubernetes、MLflow、CI/CD}
\section{榮譽}

\cventry{2021年11月}
        {亞太數據科學協會}
        {亞太數據科學競賽亞軍}
        {}
        {}
        {與團隊成員於 300 多支參賽隊伍中脫穎而出，運用機器學習技術預測城市交通流量，獲亞軍。}

\cventry{2019年7月}
        {香港大學統計學系}
        {最佳畢業論文獎}
        {}
        {}
        {畢業論文「基於深度學習的香港樓價預測模型」獲評為年度最佳碩士論文。}
\section{證書}

\cventry{2022年3月}
        {Amazon Web Services}
        {AWS Certified Machine Learning – Specialty}
        {\url{https://aws.amazon.com/certification/}}
        {}
        {}

\cventry{2021年8月}
        {Microsoft}
        {Azure Data Scientist Associate}
        {\url{https://learn.microsoft.com/en-us/credentials/}}
        {}
        {}

\cventry{2021年1月}
        {Google}
        {TensorFlow Developer Certificate}
        {\url{https://www.tensorflow.org/certificate}}
        {}
        {}
\section{出版刊物}

\cventry{2022年6月}
        {A Hybrid Deep Learning Approach for Urban Traffic Flow Prediction}
        {IEEE Access}
        {\url{https://doi.org/10.1109/ACCESS.2022.1234567}}
        {}
        {\begin{itemize}
\item 提出結合 CNN 及 LSTM 的混合深度學習模型，用於預測香港市區交通流量
\item 實驗結果顯示所提出的模型比傳統 ARIMA 減少 25\% 預測誤差
\item 文章已獲 IEEE Access 收錄，並被多位學者引用
\end{itemize}}

\cventry{2020年12月}
        {應用自然語言處理於金融文本情感分析}
        {香港數據科學期刊}
        {\url{https://www.hkdsj.org/articles/2020/12}}
        {}
        {\begin{itemize}
\item 探討如何利用 BERT 模型分析上市公司年報及新聞稿中的情感傾向
\item 結合金融領域知識詞典，提升情感分析模型於金融文本的準確度
\end{itemize}}
\section{項目}

\cventry{2021年2月–2021年6月}
        {一個基於深度學習的香港物業估價及樓價預測平台}
        {香港樓價預測系統}
        {\url{https://github.com/wingyanchan/hk-property-price}}
        {}
        {\begin{itemize}
\item 整合差餉物業估價署、土地註冊處及公共交通數據
\item 使用 XGBoost 及深度神經網絡建立樓價預測模型，R² 達 0.91
\item 開發 Flask API 及 React 前端，讓用戶輸入物業資料即時獲得估價
\end{itemize}
\textbf{關鍵字}：Python、XGBoost、React、房地產、深度學習}

\cventry{2022年9月–2023年3月}
        {為銀行客戶服務部門設計的粵語智能聊天機器人}
        {智能客戶服務聊天機器人}
        {\url{https://github.com/wingyanchan/chatbot}}
        {}
        {\begin{itemize}
\item 利用 RAG（檢索增強生成）架構結合領域知識庫，解答客戶常見問題
\item 使用 LangChain 及 OpenAI API 開發，並部署至企業 Kubernetes 集群
\item 聊天機器人成功處理 70\% 的常見查詢，大幅減輕前線客服人員負擔
\end{itemize}
\textbf{關鍵字}：RAG、LangChain、Kubernetes、生成式 AI、客戶服務}
\section{興趣愛好}

\cvline{人工智能倫理}{負責任AI、數據私隱、公平性}
\cvline{戶外活動}{行山、露營、攝影}
\section{志願服務}

\cventry{2021年1月至今}
        {兼職導師}
        {香港數據科學義工網絡}
        {\url{https://www.datasciencevolunteer.hk}}
        {}
        {\begin{itemize}
\item 為非牟利機構提供免費數據分析諮詢
\item 定期主持數據科學工作坊，教授 Python 及數據可視化技巧
\item 協助設計「社企數據儀表板」，用於追蹤服務成效
\end{itemize}}

\end{document}